\chapter{Inteligência Artificial e Educação Matemática: contexto e desafios}
\label{cap:contexto}

% =====================================================================
% PAPEL DESTE CAPÍTULO (registro pedagógico/educacional):
%   Responde "por que isso importa para o professor".
%   NÃO explica arquitetura transformer em detalhe, NÃO lista as
%   estratégias de mitigação, NÃO enuncia a hipótese arquitetural por
%   inteiro. Esses conteúdos técnicos pertencem ao Capítulo 2.
% =====================================================================

\section{Inteligência Artificial: definição operacional e classificações relevantes}
\label{sec:ia-definicao}

O termo \emph{inteligência artificial} foi empregado formalmente pela primeira vez em
1956, durante a Conferência de Dartmouth, convocada por \citet{mccarthy1955dartmouth} para
reunir pesquisadores interessados em simular o raciocínio humano por meio de máquinas. Antes
disso, Alan Turing já havia lançado, em 1950, uma questão fundadora --- \emph{``Can machines
think?''} --- e proposto um critério comportamental para avaliá-la: se uma máquina produz
respostas indistinguíveis das humanas em uma interação textual mediada exclusivamente por
escrita, considera-se que ela demonstra comportamento inteligente \citep{turing1950computing}.
Não por acaso, a primeira aplicação formal dessa ideia, o \emph{Logic Theorist}
\citep{newell1956logic}, foi construída para demonstrar teoremas da lógica proposicional ---
evidenciando, desde a origem, o vínculo estreito entre inteligência artificial e raciocínio
matemático.

Para os fins desta dissertação, entretanto, adota-se uma definição funcional em lugar de uma
genealogia histórica. Segundo \citet{russell2021aima}, um sistema é dito inteligente quando
percebe seu ambiente, processa informações e age de modo a maximizar a realização de objetivos
previamente definidos. Essa perspectiva ancora a discussão no comportamento observável dos
sistemas --- e não em questões filosóficas sobre consciência ou cognição ---, o que é suficiente
e adequado ao contexto educacional aqui investigado.

Do ponto de vista das classificações, distinguem-se duas categorias relevantes para este
trabalho. A primeira é a chamada \textbf{IA estreita} (\emph{narrow AI}), que designa sistemas
projetados para executar tarefas específicas com alto desempenho --- reconhecimento de imagens,
tradução automática, sistemas de recomendação --- sem capacidade de generalização além do
domínio para o qual foram treinados. A segunda é a \textbf{IA generativa}, categoria mais
recente, que compreende modelos capazes de produzir conteúdo original --- texto, imagem, código
--- a partir de padrões aprendidos em grandes volumes de dados. É nessa segunda categoria que se
inserem os sistemas centrais desta dissertação.

Para os propósitos deste trabalho, entende-se por ``sistema de IA'' qualquer ferramenta
computacional baseada em modelos de linguagem de grande escala (\emph{Large Language Models} ---
LLMs) capaz de gerar, avaliar ou organizar conteúdo textual de natureza matemática. Essa
delimitação exclui, deliberadamente, sistemas voltados à visão computacional, robótica ou jogos,
cujas arquiteturas e desafios diferem substancialmente do escopo aqui investigado.

\section{IA generativa e modelos de linguagem de grande escala}
\label{sec:ia-generativa}

% NOTA DE REVISÃO: esta seção foi enxugada. A explicação detalhada da
% arquitetura transformer, da inferência probabilística e do fato de que
% "o modelo não computa, prevê" foi REMOVIDA daqui e vive agora na
% Seção \ref{sec:llm} (Capítulo 2), sua casa canônica. Aqui fica apenas
% o panorama histórico que situa o leitor.

A trajetória que levou aos atuais modelos de linguagem de grande escala tem início nas primeiras
tentativas de simular a comunicação humana por meio de máquinas. Na década de 1960,
\citet{weizenbaum1966eliza} desenvolveu o ELIZA, programa estruturado para estabelecer diálogos
por meio da identificação de padrões linguísticos e da reorganização das frases apresentadas pelo
usuário. Ainda que não houvesse compreensão real do conteúdo, o ELIZA demonstrou que máquinas
poderiam sustentar interações textuais convincentes e abriu caminho para o desenvolvimento de
tecnologias conversacionais progressivamente mais sofisticadas.

Nas décadas seguintes, o avanço do poder computacional, a ampliação das bases de dados
disponíveis e o desenvolvimento de métodos conexionistas --- especialmente as redes neurais
artificiais --- reposicionaram a inteligência artificial como campo estratégico de pesquisa.
Nesse contexto, \citet{bengio2003neural} propuseram os primeiros modelos neurais de linguagem
capazes de representar palavras como vetores em espaços semânticos contínuos, estabelecendo as
bases conceituais para o que viria a seguir. O avanço decisivo ocorreu em 2017, com a introdução
da arquitetura \emph{transformer} por \citet{vaswani2017attention}, que se tornou a base
tecnológica da geração contemporânea de LLMs.

Os \textbf{modelos de linguagem de grande escala} são a realização desse percurso em escala
massiva: modelos como o GPT-3 \citep{brown2020gpt3}, treinados com centenas de bilhões de
parâmetros, demonstraram que o aumento de escala produz capacidades não observadas em modelos
menores --- geração de código, tradução entre idiomas e produção de texto de alta qualidade ---,
o que os tornou a tecnologia dominante para geração de conteúdo textual na presente década. Para
os fins desta dissertação, importa reter uma propriedade central desses modelos: sua geração é
\textbf{probabilística e não-determinística}. Essa característica confere grande versatilidade ---
inclusive para gerar questões matemáticas originais ---, mas também os torna suscetíveis a erros
em tarefas que exigem rigor lógico e consistência simbólica. O funcionamento técnico que explica
essa propriedade --- e a distinção fundamental entre \emph{prever} e \emph{computar} --- é
desenvolvido na Seção~\ref{sec:llm}.

\section{IA na Educação: panorama de aplicações}
\label{sec:ia-educacao}

A incorporação da inteligência artificial ao campo educacional configura-se como um processo mais
complexo e gradual do que o observado em outros setores. Diferentemente do ambiente empresarial,
cuja adoção é fortemente impulsionada por demandas de produtividade e competitividade, a educação
envolve dimensões pedagógicas, éticas e sociais que exigem reflexão crítica e regulamentação
cuidadosa. Questões relacionadas à formação docente, à infraestrutura tecnológica, à proteção de
dados de estudantes e às implicações didáticas do uso dessas ferramentas tornam o avanço
necessariamente mais cauteloso. Assim, embora a IA apresente potencial significativo para apoiar
práticas pedagógicas e personalizar a aprendizagem, sua incorporação ao contexto escolar ainda
ocorre de maneira progressiva e permeada por debates acerca de seus limites e possibilidades.

Dentre as categorias de aplicação mais estudadas na literatura, os \textbf{sistemas tutoriais
inteligentes} (\emph{Intelligent Tutoring Systems} --- ITS) figuram entre as mais antigas e
consolidadas. Esses sistemas buscam simular a interação individualizada entre professor e aluno,
identificando lacunas no conhecimento do estudante e adaptando a sequência de atividades em tempo
real. A motivação teórica remonta ao estudo seminal de \citet{bloom1984twosigma}, que demonstrou
que o atendimento individualizado produz ganhos de aprendizagem significativamente superiores aos
do ensino coletivo convencional --- o chamado \emph{problema das duas sigmas}. Pesquisas
posteriores indicam que ITS bem projetados conseguem aproximar-se desse referencial, com efeitos
positivos documentados especialmente em matemática e ciências.

Uma segunda categoria relevante compreende os \textbf{sistemas de recomendação e ambientes
adaptativos}. Nesses sistemas, algoritmos de aprendizado de máquina analisam o histórico de
interações do estudante --- desempenho em atividades, tempo de resposta, padrões de erro --- e
recomendam conteúdos, exercícios ou percursos de aprendizagem personalizados. A lógica subjacente
é semelhante à dos sistemas de recomendação comerciais, mas orientada a objetivos pedagógicos:
maximizar a progressão do estudante respeitando seu ritmo e suas dificuldades específicas.
Plataformas como Khan Academy e Duolingo incorporam variações dessa abordagem em larga escala.

A \textbf{geração automática de conteúdo educacional} constitui a terceira categoria e a mais
diretamente relacionada ao objeto desta dissertação. Com o advento dos LLMs, tornou-se
tecnicamente viável produzir questões, enunciados, resumos, \emph{feedback} personalizado e até
planos de aula de forma automatizada, a partir de especificações fornecidas pelo professor. Essa
possibilidade tem despertado interesse crescente na pesquisa em Educação Matemática, pois pode
reduzir o tempo dedicado à elaboração de materiais e ampliar a diversidade de situações-problema
disponíveis para o trabalho em sala de aula. Entretanto, como discutido na
Seção~\ref{sec:ia-matematica}, a geração automática em contextos matemáticos apresenta
fragilidades específicas que demandam estratégias de verificação e validação docente.

Por fim, os \textbf{assistentes conversacionais} representam a categoria de maior penetração no
cotidiano de professores e estudantes na atualidade. Ferramentas baseadas em LLMs --- como o
ChatGPT, o Gemini e o próprio Claude --- são amplamente utilizadas para tirar dúvidas, solicitar
explicações alternativas, gerar exemplos e apoiar a resolução de problemas. Sua acessibilidade,
disponibilidade contínua e capacidade de adaptar o nível de linguagem ao interlocutor conferem-lhes
potencial pedagógico relevante. Ao mesmo tempo, sua natureza probabilística --- que pode produzir
respostas plausíveis, porém incorretas, sem qualquer sinalização de incerteza --- impõe ao
professor o papel insubstituível de mediador crítico, tema retomado ao final deste capítulo.

\section{IA no ensino de Matemática: potencialidades e fragilidades}
\label{sec:ia-matematica}

A chegada de modelos de linguagem de grande escala ao cotidiano escolar suscita tanto entusiasmo
quanto cautela entre educadores e pesquisadores. No campo da Matemática, essa ambivalência é
especialmente marcada: ao mesmo tempo que os LLMs demonstram capacidade notável de comunicar
ideias matemáticas em linguagem natural e de produzir resoluções passo a passo, eles exibem
falhas sistemáticas em tarefas que deveriam ser triviais para qualquer computador convencional.

\subsection*{Potencialidades reconhecidas}

Do ponto de vista pedagógico, os LLMs oferecem três contribuições que merecem destaque. A
primeira é a \textbf{multiplicidade de abordagens explicativas}: diante de uma mesma questão, o
modelo pode reformular o enunciado, apresentar exemplos análogos, oferecer representações visuais
em ASCII ou sugerir estratégias heurísticas alternativas, adaptando a linguagem ao perfil do
estudante. Essa flexibilidade aproxima-se do ideal de ensino diferenciado preconizado pela
literatura em educação matemática.

A segunda potencialidade é a \textbf{geração rápida de material didático}. Um professor pode, em
poucos minutos, obter rascunhos de enunciados, variações de um mesmo problema ou conjuntos de
exercícios graduados por dificuldade, o que reduz a carga de criação de conteúdo e amplia a
personalização do ensino.

A terceira contribuição é o \textbf{suporte à articulação verbal do raciocínio}. Ao se pedir que o
modelo apresente cada etapa de resolução, o processo de solução torna-se legível e discutível, o
que pode ser aproveitado pelo professor como material para análise e debate em sala. O fundamento
técnico desse recurso --- a técnica de \emph{chain-of-thought prompting} --- e seus limites são
tratados na Seção~\ref{sec:mitigacao}.

\subsection*{Fragilidades documentadas}

As limitações dos LLMs em Matemática são igualmente documentadas e, para os fins desta
dissertação, mais relevantes do que as potencialidades.
\citet{frieder2023mathematical} conduziram uma avaliação sistemática das capacidades matemáticas
do ChatGPT/GPT-4, concluindo que o modelo apresenta desempenho aceitável em questões de rotina,
mas falha de forma expressiva em problemas que exigem demonstração formal, manipulação simbólica
complexa ou raciocínio em múltiplos passos encadeados. Essa assimetria entre \emph{fluência
discursiva} e \emph{rigor computacional} é a característica mais decisiva para o uso educacional da
tecnologia: o modelo escreve sobre Matemática com desenvoltura, mas não garante a correção daquilo
que escreve.

Do ponto de vista do professor, o que importa reter neste capítulo é a consequência prática dessa
assimetria --- erros de aritmética, de álgebra simbólica e de encadeamento lógico ocorrem de forma
recorrente e pouco previsível. A tipologia técnica desses erros, com seus mecanismos e a
literatura que os documenta, é apresentada na Seção~\ref{sec:limitacoes}.

\subsection*{O fenômeno do erro silencioso}
\label{sec:erro-silencioso}

% CASA CANÔNICA do conceito de "erro silencioso": é o gancho que motiva
% todo o trabalho, e por isso é desenvolvido por extenso AQUI. Nos demais
% capítulos ele deve ser apenas referenciado (\ref{sec:erro-silencioso}).

A principal implicação pedagógica das fragilidades acima não é a possibilidade de erro em si, mas o
modo como esse erro se apresenta: o LLM tende a produzir respostas \textbf{plausíveis e
incorretas}, redigidas com o mesmo tom de confiança das respostas corretas. Diferentemente de uma
calculadora, que simplesmente rejeita uma entrada inválida, ou de um sistema de computação
algébrica, que emite uma mensagem de erro, o LLM elabora uma explicação coerente para um resultado
equivocado. Esse fenômeno --- denominado \emph{alucinação} na literatura de processamento de
linguagem natural e aqui chamado de \textbf{erro silencioso} por sua invisibilidade ao leitor não
especializado --- representa um risco específico no contexto educacional, em que o professor ou o
estudante pode aceitar um gabarito incorreto sem perceber o equívoco.

\subsection*{Implicações para esta dissertação}

O diagnóstico apresentado neste capítulo fundamenta a hipótese central deste trabalho: o uso
isolado de um LLM para geração de questões matemáticas é insuficiente para garantir correção e
qualidade didática. A resposta a esse diagnóstico --- combinar o poder expressivo do LLM com uma
verificação determinística e com uma avaliação pedagógica estruturada --- não é um refinamento
opcional, mas uma necessidade arquitetural. Os fundamentos técnicos dessa combinação, e o
argumento completo que a sustenta, são desenvolvidos no Capítulo~\ref{cap:tecnica}.
